%\documentclass[12pt,a4paper]{book}
\documentclass[a4paper,french]{book}
\usepackage[17pt]{extsizes}


\input{../1ere-preambule}

\newcommand{\va}[1]{\lvert \ #1 \ \rvert}
\newcommand{\vaf}[1]{\left\lvert \ #1\  \right\rvert}

\begin{document}
\graphicspath{{../images/}}
\setcounter{chapter}{1}
\chapter[Correction des exercices]{Trigonométrie \\ correction des exercices}
\thispagestyle{fancy}

\section*{Mesurer un angle en radian}

\begin{multicols}{2}
\subsection*{Ex 16 p 193}

\begin{enumMet}
    \item triangle équilateral : 3 angles de $\dfrac{\pi}{3}$ rad
    \item triangle rectangle isocèle : 1 angle de $\dfrac{\pi}{2}$ rad. 2 angles de $\dfrac{\pi}{4}$ rad
\end{enumMet}

\columnbreak
\subsection*{Ex 18 p 193}

\begin{itemMet}
    \item $\dfrac{\pi}{7}$ : $\dfrac{\pi}{7}+2\pi=\dfrac{\pi}{7}+\dfrac{14\pi}{7}=\dfrac{15\pi}{7}$
    \item $-\dfrac{\pi}{5}$ : $-\dfrac{\pi}{5}+2\pi=-\dfrac{\pi}{5}+\dfrac{10\pi}{5}=\dfrac{9\pi}{5}$
    \item $\dfrac{\pi}{6}$ : $\dfrac{\pi}{6}+2\pi=-\dfrac{\pi}{6}+\dfrac{12\pi}{6}=\dfrac{13\pi}{6}$
    \item $\dfrac{2\pi}{11}$ : $\dfrac{2\pi}{11}+2\pi=\dfrac{2\pi}{11}+\dfrac{22\pi}{11}=\dfrac{24\pi}{11}$
    \item $-\dfrac{7\pi}{4}$ : $-\dfrac{7\pi}{4}+2\pi=-\dfrac{7\pi}{4}+\dfrac{8\pi}{4}=\dfrac{\pi}{4}$
\end{itemMet}

\end{multicols}

 \begin{multicols}{2}
\subsection*{Ex 20 p 193}

\includegraphics[scale=0.7]{./corr-ch02/1.02.02-01.png}
\columnbreak
\subsection*{Ex 21 p 193}
$\dfrac{3 \pi}{2}$ est associé au même point que $\dfrac{-\pi}{2}$ et $\dfrac{7 \pi}{4}$ est associé au même point que $\dfrac{-\pi}{4}$. On obtient la figure ci-dessous.\\
\includegraphics[scale=0.7]{./corr-ch02/1.02.02-02.png}
 \end{multicols}
\pagebreak


 \subsection*{Ex 23 p 193}

 On a $\dfrac{13\pi}{6}=\dfrac{\pi}{6}+2\pi$, $-\dfrac{11\pi}{6}=\dfrac{\pi}{6}-2\pi$ et $\dfrac{25\pi}{6}=\dfrac{\pi}{6}+4\pi$ qui représentent le même point correspondant à $\dfrac{\pi}{6}$.\\
 On a également $\dfrac{11\pi}{6}=-\dfrac{\pi}{6}+2\pi$ et $-\dfrac{13\pi}{6}=-\dfrac{\pi}{6}-2\pi$ qui représentent le même point correspondant à $-\dfrac{\pi}{6}$.\\

\subsection*{Ex 24 p 193}
\begin{multicols}{2}
    On utilise le fait que les radians et les degrés sont proportionnels et le fait que  $\pi$ rad soit égal à $180\degree$ :
    \begin{enumMet}
        \item $\dfrac{\pi}{16}\text{\ rad}=\dfrac{180\degree}{16}=11,25\degree$
        \item $\dfrac{\pi}{72}\text{\ rad}=\dfrac{180\degree}{72}=2,5\degree$
        \item $\dfrac{\pi}{7}\text{\ rad}=\dfrac{180\degree}{7}\approx 25,7\degree$
    \end{enumMet}
%\end{multicols}

\subsection*{Ex 25 p 193}
On utilise la formule suivante : radian $=\dfrac{\text { degre } \times \pi}{180}$
    \begin{enumMet}
        \item $\dfrac{75 \times \pi}{180}=\dfrac{5 \pi}{12} \mathrm{rad}$
        \item $\dfrac{120 \times \pi}{180}=\dfrac{2 \pi}{3} \mathrm{rad}$
        \item $\dfrac{405 \times \pi}{180}=\dfrac{9 \pi}{4} \mathrm{rad}$
    \end{enumMet}
\end{multicols}

 \subsection*{Ex 41 p 194}

\begin{enumMet}
    \item $\dfrac{2\pi}{3}$ se trouvera dans le quadrant 2 car $\dfrac{\pi}{2}<\dfrac{2\pi}{3}<\pi$
    \item $\dfrac{\pi}{8}$ se trouvera dans le quadrant 1 car $0<\dfrac{\pi}{8}<\dfrac{\pi}{2}$
    \item $\dfrac{5\pi}{4}$ se trouvera dans le quadrant 3 car $\pi<\dfrac{5\pi}{4}<\dfrac{3\pi}{2}$
    \item $-\dfrac{2\pi}{3}$ se trouvera dans le quadrant 3 car $-\pi<-\dfrac{2\pi}{3}<-\dfrac{\pi}{2}$
\end{enumMet}

\vspace{-0.2cm}

\subsection*{Ex 43 p 194}

\begin{enumMet}
    \item Le nombre réel $\dfrac{-4\pi}{3}$ correspond au point $D$.
    \item Le nombre réel $\dfrac{7\pi}{2}$ correspond au point $J'$.
    \item Le nombre réel $\dfrac{-11\pi}{6}$ correspond au point $A$.
    \item Le nombre réel $13\pi$ correspond au point $I'$.
\end{enumMet}

\vspace{-0.2cm}

\subsection*{Ex 44 p 194}

\begin{enumMet}
    \item $2018\pi=2\pi\times1009$.Cela revient donc à faire 1009 tours entiers du cercle trigonométrique. Ce réel correspond donc au point $I$. 
    \item $\dfrac{403\pi}{2}=2\pi\times100+\dfrac{3\pi}{2}$. Cela revient donc à faire 100 tours entiers du cercle trigonométrique dans le sens trigonométrique puis les trois quarts d’un autre tour. Ce réel correspond donc au point $J'$.
    \item $-\dfrac{47\pi}{3}=-2\pi\times8+\dfrac{\pi}{3}$. Cela revient donc à faire 16 tours entiers du cercle trigonométrique dans le sens indirect puis le tier d’un autre demi-tour dans le sens trigonométrique. Ce réel correspond donc au point $C$.
    \item $\dfrac{109\pi}{4}=2\pi\times13+\dfrac{5\pi}{4}$. Cela revient donc à faire 16 tours entiers du cercle trigonométrique dans le sens trigonométrique puis un demi-tour suivi d'un quart de demi-tour dans le sens trigonométrique. Ce réel correspond donc au point $H$.
\end{enumMet}

\subsection*{Ex 46 p 195}

\begin{enumMet}
    \item $\dfrac{3\pi}{2}=-\dfrac{\pi}{2}+2\pi=-\dfrac{5\pi}{2}+2\times2\pi$. L'intrus est donc $\dfrac{9\pi}{2}$
    \item $\dfrac{14\pi}{3}=-\dfrac{8\pi}{6}+3\times2\pi=-\dfrac{10\pi}{3}+4\times2\pi$. L'intrus est donc $\dfrac{\pi}{3}$
    \item $\dfrac{7\pi}{4}=-\dfrac{\pi}{4}+2\pi=-\dfrac{9\pi}{4}+2\times2\pi$. L'intrus est donc $-\dfrac{19\pi}{4}$
    \item $\sqrt{9}\pi=3\pi=\pi+2\pi=-\pi+2\times2\pi$. L'intrus est donc $0$
\end{enumMet}

\subsection*{Ex 49 p 195}

\begin{itemMet}
    \item Comme le triangle $FGH$  est isocèle en $F$ , ses angles à la base sont de même mesure d’où : \\
$$\widehat{FHG}=\dfrac{\pi-\dfrac{\pi}{5}}{2}=\dfrac{2\pi}{5}$$
    \item Comme le triangle $BCE$ est équilatéral, $\widehat{BEC}=\dfrac{\pi}{3}$ 
    \item Comme $ABCD$ est un carré, $\widehat{ABC}=\dfrac{\pi}{2}$ rad et donc :\\ $\widehat{ABE}=\widehat{ABC}-\widehat{EBC}=\dfrac{\pi}{2}-\dfrac{\pi}{3}=\dfrac{\pi}{6}$ rad.
\end{itemMet}

\vspace*{-.2cm}

\subsection*{Ex 51 p 195}
On utilise la proportionnalité des mesures d'angles et on sait que $360\degree=2\pi$ rad.\\
\begin{tabular}{ | Xc | Xc | Xc | Xc | Xc | Xc |}
    \hline			
    Mesure d'angle en degrée($\degree$)  & 20 & 168 & 26 & 14 & 245 \\ \hline
    Mesure d'angle en radian & $\dfrac{\pi}{9}$ & $\dfrac{14\pi}{15}$ & $\dfrac{\pi}{7}$ & $\dfrac{\pi}{13}$ & $\dfrac{49\pi}{36}$ \\
    \hline  
\end{tabular}

\subsection*{Ex 55 p 196}

\begin{enumMet}
	\item Entre $12 \mathrm{~h}$ et $12 \mathrm{~h} 20$ l'aiguille a parcourue une distance de $\dfrac{2 \pi}{3}$ unités.
	\item Entre 15 h15 et $16 \mathrm{~h} 30$ l'aiguille a parcourue une distance de $\dfrac{5 \pi}{2}$ unités.
	\item Entre $20 \mathrm{~h} 30$ et $22 \mathrm{~h} 50$ l'aiguille a parcourue une distance de $\pi+2 \pi+\dfrac{5 \pi}{3}=\dfrac{14 \pi}{3}$.
	\item Entre $14 \mathrm{~h} 50$ et $17 \mathrm{~h} 22$ l'aiguille a parcourue une distance de $2 \times \dfrac{\pi}{6}+2 \pi+2 \pi+\dfrac{22 \pi}{30}=\dfrac{76 \pi}{15} \text { unités. }$
\end{enumMet}



\subsection*{Ex 56 p 196}

\begin{enumMet}
	\item Si la roue tourne d'un angle de $\dfrac{\pi}{2}$, la flèche tombe sur le 3.
	\item Si la roue tourne d'un angle de $\dfrac{19 \pi}{10}$, la flèche tombe sur le 10 .
	\item Si la roue tourne d'un angle de $\dfrac{-5 \pi}{4}$, la flèche tombe sur le 4 .
	\item Si la roue tourne d'un angle de $\dfrac{34 \pi}{3}$, la flèche tombe sur le 7 .
\end{enumMet}


\section*{cosinus et sinus}

\subsection*{Ex 17 p 193}

\begin{tabular}{ | Xc | Xc | Xc | Xc | Xc | Xc |}
    \hline
    \   & $\dfrac{\pi}{3}$ & $-\dfrac{\pi}{3}$ & $-\dfrac{2\pi}{3}$ & $\dfrac{3\pi}{3}$ & $\dfrac{27\pi}{3}$ \\ \hline
    Cosinus  & $+\dfrac{1}{2}$ & $+\dfrac{1}{2}$ & $-\dfrac{1}{2}$ & $-1$ & $-1$ \\ \hline
    Sinus & $+\dfrac{\sqrt{3}}{2}$ & $-\dfrac{\sqrt{3}}{2}$ & $-\dfrac{\sqrt{3}}{2}$ & $0$ & $0$ \\
    \hline  
\end{tabular}



\subsection*{Ex 19 p 193}
Sur $\left[0;\dfrac{\pi}{2}\right]$, on a :
\begin{enumMet}
    \item $\cos(x)=0,5\iff x=\dfrac{\pi}{3}$
    \item $\sin(x)=0,5\iff x=\dfrac{\pi}{6}$
    \item $\sin(x)=\cos(x)\iff x=\dfrac{\pi}{4}$
\end{enumMet}

\subsection*{Ex 59 p 196}

\begin{tabular}{ | Xc | Xc | Xc | Xc | Xc |}
    \hline
    M est dans le quadrant & $1$ & $2$ & $3$ & $4$ \\ \hline
    signe de $\cos(x)$  & $+$ & $-$ & $-$ & $+$ \\ \hline
    signe de $\sin(x)$ & $+$ & $+$ & $-$ & $-$ \\
    \hline  
\end{tabular}

\subsection*{Ex 60 p 196}

\begin{tabular}{ | Xc | Xc | Xc | Xc | Xc | Xc |}
    \hline
    $x$  & $\dfrac{2\pi}{3}$ & $\dfrac{3\pi}{4}$ & $-\dfrac{5\pi}{6}$ & $-\dfrac{3\pi}{4}$ & $-\dfrac{\pi}{2}$ \\ \hline
    Point image  & $D$ & $E$ & $G$ & $H$ & $I'$ \\ \hline
    $\cos(x)$  & $-\dfrac{1}{2}$ & $-\dfrac{\sqrt{2}}{2}$ & $-\dfrac{\sqrt{3}}{2}$ & $-\dfrac{\sqrt{2}}{2}$ & $0$ \\ \hline
    $\sin(x)$ & $\dfrac{\sqrt{3}}{2}$ & $\dfrac{\sqrt{2}}{2}$ & $-\dfrac{1}{2}$ & $-\dfrac{\sqrt{2}}{2}$ & $-1$ \\
    \hline  
\end{tabular}

\subsection*{Ex 61 p 197}

\begin{tabular}{ | Xc | Xc | Xc | Xc | Xc | Xc |}
    \hline
    $x$  & $\dfrac{5\pi}{6}$ & $-\pi$ & $-\dfrac{2\pi}{3}$ & $-\dfrac{\pi}{4}$ & $-\dfrac{\pi}{6}$ \\ \hline
    Point image  & $F$ & $J'$ & $K$ & $M$ & $N$ \\ \hline
    $\cos(x)$  & $-\dfrac{\sqrt{3}}{2}$ & $-1$ & $-\dfrac{1}{2}$ & $\dfrac{\sqrt{2}}{2}$ & $\dfrac{\sqrt{3}}{2}$ \\ \hline
    $\sin(x)$ & $\dfrac{1}{2}$ & $0$ & $-\dfrac{\sqrt{3}}{2}$ & $-\dfrac{\sqrt{2}}{2}$ & $-\dfrac{1}{2}$ \\
    \hline  
\end{tabular}

\vspace{-.2cm}

\subsection*{Ex 66 p 197}
On cherche les valeurs correspondantes à l'aide des valeurs connues et de la position sur le cercle trigonométrique.
\begin{enumMet}
    \item
    \begin{enumMet}
        \item Sur $\left[0;\pi\right]$ : $\cos(x)=\dfrac{-\sqrt{2}}{2}\iff x=\dfrac{3\pi}{4}$
        \item Sur $\left[-\pi;-\dfrac{\pi}{2}\right]$ : $\cos(x)=\dfrac{-\sqrt{2}}{2}\iff x=-\dfrac{3\pi}{4}$
    \end{enumMet}
    \item
    \begin{enumMet}
        \item Sur $\left[\dfrac{\pi}{2};\pi\right]$ : $\sin(x)=\dfrac{1}{2}\iff x=\dfrac{5\pi}{6}$
        \item Sur $\left[-\dfrac{\pi}{2};\dfrac{\pi}{2}\right]$ : $\sin(x)=\dfrac{1}{2}\iff x=-\dfrac{\pi}{6}$
    \end{enumMet}
\end{enumMet}

\subsection*{Ex 67 p 197}
\begin{enumMet}
    \item
    \begin{enumMet}
        \item Sur $\left[\pi;2\pi\right]$ : $\cos(x)=\dfrac{\sqrt{3}}{2}\iff x=\dfrac{11\pi}{6}$
        \item Sur $\left[\dfrac{\pi}{2};\dfrac{3\pi}{2}\right]$ : $\cos(x)=\dfrac{\sqrt{3}}{2}\iff S=\varnothing$
    \end{enumMet}
    \item
    \begin{enumMet}
        \item Sur $\left[-\pi;\pi\right]$ : $\sin(x)=-1\iff x=\dfrac{-\pi}{2}$
        \item Sur $\left[\dfrac{35\pi}{2};\dfrac{37\pi}{2}\right]$ : $\sin(x)=-1\iff x=\dfrac{35\pi}{2}$
    \end{enumMet}
    \item Sur $\left[0;2\pi\right]$ : $\cos(x)=3\iff S=\varnothing$ car pour tout $x\in \R$, $-1\le \cos(x) \le 1$
\end{enumMet}

\end{document}

